\section{Limitations}

This framework reimplements a thin slice of what LLVM's \texttt{instcombine} and
dead code passes already do, and it does so on purpose, to keep every rewrite
small enough to verify exhaustively. It is not a general optimizer and does not
try to be.

The scope is scalar integer IR. Floating point, vectors, volatile, and atomic
operations are recognized and skipped rather than handled. Common subexpression
elimination is block local, so redundancy that spans blocks, which a global value
numbering pass would catch, is left in place. The transforms do not fold branches
or simplify the control flow graph, which is why a real \texttt{opt -O1} pulls
ahead on the control flow heavy kernels.

The timing results are from a single machine running under WSL2, where there is
no control over the CPU frequency governor, so a few milliseconds of jitter are
expected. The kernels run in tens of milliseconds, so for the kernels with a
small structural change the runtime difference falls within the measurement
spread and is reported as structural only rather than dressed up as a speedup.
The one kernel with substantial redundancy shows a runtime improvement clearly
outside that spread.
